\begin{figure}[h!]
\begin{center}
\includegraphics[width=0.60\columnwidth]{figures/ID3183_class_3_target_4.png}
\end{center}
\caption{ Classification: 3 actual target: 4}
\label{fig:ID3183_class_3_target_4}
\end{figure}
